\chapter{Périmètre et aire}\label{ch:g5:areas}

L'année~4 comptait les carreaux du quadrillage
(\cref{ch:g4:measure})~; ce chapitre gagne la première
\emph{formule} d'aire --- longueur fois largeur pour le rectangle ---
et apprend les unités cm$^2$ et m$^2$. Les formules pour les
triangles et d'autres formes mûrissent dans \cref{ch:g6:measure} et
\cref{ch:g7:areas}.

\section{Deux mesures différentes}

\begin{example}[Bord contre surface]\label{ex:g5:areas:contrast}
Un jardinier a besoin du \emph{périmètre} pour acheter la clôture et
de l'\emph{aire} pour acheter le gazon. Les deux ne se suivent
pas~: étirer un rectangle plus fin et plus long peut garder son
aire tandis que son périmètre grandit --- comparer un
$6 \times 2$ et un $4 \times 3$ (même aire $12$, périmètres $16$ et
$14$).
\end{example}

\section{La formule du rectangle}

\begin{proposition}[Aire d'un rectangle]\label{prop:g5:areas:rectangle}
Un rectangle de longueur $L$ et de largeur $\ell$ (dans la même
unité) a pour aire
\[
A = L \times \ell ,
\]
en unités d'aire~: cm$^2$ (carrés de côté $1$~cm), m$^2$ (carrés de
côté $1$~m), \dots
\end{proposition}

\begin{proof}[Pourquoi, par comptage]
Recouvrir le rectangle de carreaux unitaires~: $\ell$ rangées de
$L$ carreaux chacune, donc $L \times \ell$ carreaux en tout ---
exactement le rectangle de la multiplication de \cref{ch:g3:mult}.
\end{proof}

\begin{omfigure}
\begin{tikzpicture}[scale=0.62]
  \draw[gray!50, thin] (0,0) grid (6,3);
  \draw[very thick, omDef] (0,0) rectangle (6,3);
  \fill[omThm!30] (0,2) rectangle (1,3);
  \draw[omThm, thick] (0,2) rectangle (1,3);
  \node[right, font=\small, omThm] at (1.15,2.5) {$1$ cm$^2$};
  \node[below, font=\small] at (3,-0.15) {$6$ cm};
  \node[left, font=\small] at (0,1.5) {$3$ cm};
  \node[font=\small] at (3.9,1.2) {$6 \times 3 = 18$ cm$^2$};
\end{tikzpicture}

{\small Trois rangées de six centimètres-carrés~: la formule
$L \times \ell$ est le comptage, écrit une fois pour toutes.}
\end{omfigure}

\begin{example}\label{ex:g5:areas:rectangle}
Un tapis mesure $2.5$~m sur $2$~m~: aire $2.5 \times 2 = 5$~m$^2$.
Un timbre mesure $3$~cm sur $2.4$~cm~: aire $3 \times 2.4 = 7.2$
cm$^2$. Même formule, n'importe quelle unité --- tant que les deux
côtés utilisent la \emph{même}.
\end{example}

\begin{example}[Les unités d'aire se convertissent par 100]\label{ex:g5:areas:units}
$1$~m $= 100$~cm, mais $1$~m$^2 = 100 \times 100 = 10\,000$~cm$^2$~:
un mètre carré est une grille $100$ par $100$ de centimètres
carrés. Les unités d'aire sautent par centaines, pas par dizaines.
\end{example}

\section{Figures composées}

\begin{method}[Couper, additionner, soustraire]\label{met:g5:areas:composite}
Pour une figure faite de rectangles (un L, un T, un cadre)~:
\begin{enumerate}
  \item la découper en rectangles, ou la compléter en un grand
        rectangle~;
  \item calculer chaque aire rectangulaire avec la formule~;
  \item additionner les morceaux --- ou soustraire le trou du grand
        rectangle~;
  \item vérifier par un comptage approximatif de carreaux.
\end{enumerate}
\end{method}

\begin{example}\label{ex:g5:areas:composite}
Une pièce en L~: un rectangle $6$~m $\times$ $4$~m avec un coin
$2$~m $\times$ $2$~m manquant.
\[
\text{Aire} = 6 \times 4 - 2 \times 2 = 24 - 4 = 20 \text{ m}^2 .
\]
Ou découper le L en une bande $6 \times 2$ et une bande
$4 \times 2$~: $12 + 8 = 20$~m$^2$ --- deux chemins, une réponse.
\end{example}

\begin{example}[La moitié d'un rectangle]\label{ex:g5:areas:halfrect}
Couper un rectangle le long d'une diagonale donne deux triangles de
même aire~: chacun est la \emph{moitié} du rectangle. Un triangle
rectangle de côtés de l'angle droit $6$~cm et $4$~cm a donc pour
aire $\frac{6 \times 4}{2} = 12$~cm$^2$ --- une image à retenir
pour \cref{ch:g6:measure}, où elle devient une formule.
\end{example}

\begin{omfigure}
\begin{tikzpicture}[scale=0.62]
  \draw[gray!50, thin] (0,0) grid (6,4);
  \draw[very thick] (0,0) rectangle (6,4);
  \fill[omDef!25] (0,0) -- (6,0) -- (0,4) -- cycle;
  \draw[very thick, omDef] (6,0) -- (0,4);
  \node[font=\small] at (1.7,1.2) {moitié};
  \node[font=\small] at (4.3,2.8) {moitié};
\end{tikzpicture}

{\small La diagonale coupe le rectangle $6 \times 4$ en deux
triangles égaux de $12$ carreaux chacun.}
\end{omfigure}

\section{Exercices}

\begin{exercise}[$\star$]\label{exo:g5:areas:1}
Calculer l'aire et le périmètre d'un rectangle $8$~cm $\times$
$5$~cm. Quelle réponse est en cm, laquelle en cm$^2$~?
\end{exercise}

\begin{exercise}[$\star$]\label{exo:g5:areas:2}
Calculer les aires~: un carré de côté $9$~cm~; un rectangle
$12$~m $\times$ $7$~m~; un rectangle $4.5$~cm $\times$ $6$~cm.
\end{exercise}

\begin{exercise}[$\star$]\label{exo:g5:areas:3}
Un rectangle a une aire de $63$~cm$^2$ et une longueur de $9$~cm.
Trouver sa largeur, puis son périmètre.
\end{exercise}

\begin{exercise}[$\star$]\label{exo:g5:areas:4}
Dessiner deux rectangles différents d'aire $24$ carreaux sur papier
quadrillé, et calculer les deux périmètres. Même aire --- même
périmètre~?
\end{exercise}

\begin{exercise}[$\star$]\label{exo:g5:areas:5}
Convertir~: $3$~m$^2$ en cm$^2$~; \quad $50\,000$~cm$^2$ en m$^2$.
(Rappeler \cref{ex:g5:areas:units}~: par centaines~!)
\end{exercise}

\begin{exercise}[$\star$]\label{exo:g5:areas:6}
Une figure en T est faite d'une barre horizontale $8 \times 2$ sur
une barre verticale $2 \times 5$ (en cm). Calculer son aire en
additionnant deux rectangles.
\end{exercise}

\begin{exercise}[$\star$]\label{exo:g5:areas:7}
Un cadre photo~: un rectangle $30$~cm $\times$ $20$~cm avec une
fenêtre rectangulaire $24$~cm $\times$ $14$~cm découpée. Quelle
aire de bois le cadre utilise-t-il~?
\end{exercise}

\begin{exercise}[$\star$]\label{exo:g5:areas:8}
Calculer l'aire d'un triangle rectangle de côtés de l'angle droit
$8$~cm et $6$~cm (moitié d'un rectangle,
\cref{ex:g5:areas:halfrect}).
\end{exercise}

\begin{exercise}[$\star$]\label{exo:g5:areas:9}
Un potager rectangulaire mesure $7$~m sur $4$~m. Les graines
coûtent $2$ par mètre carré, et la clôture $3$ par mètre. Calculer
le coût des graines, puis celui de la clôture.
\end{exercise}

\begin{exercise}[$\star\star$]\label{exo:g5:areas:10}
Le sol d'un couloir mesure $12$~m de long et $2$~m de large, et
doit être recouvert de carreaux carrés de côté $50$~cm. Combien de
carreaux faut-il~? (Convertir d'abord, ou compter les carreaux dans
chaque direction.)
\end{exercise}

\begin{exercise}[$\star\star$]\label{exo:g5:areas:11}
Doubler les côtés d'un rectangle $4$~cm $\times$ $3$~cm. Que
devient son périmètre~? Son aire~? (Calculer les deux avant et
après --- les deux réponses diffèrent~!)
\end{exercise}
